% ==================================================================
% Problem definition: MRCPSP -> Cheng et al. (2015) -> multi-project tool
% ramp-up with replenishable non-renewable resources. \input into the main file.
% Notation of (M3) follows code/mp_monolithic.py (periods 0..H-1, consumption
% at task start, end-of-period inventory, orders arrive L_m periods later).
% ==================================================================
\providecommand{\calP}{\mathcal{P}}
\providecommand{\calJ}{\mathcal{J}}
\providecommand{\calQ}{\mathcal{Q}}
\providecommand{\calR}{\mathcal{R}}
\providecommand{\calM}{\mathcal{M}}
\providecommand{\calN}{\mathcal{N}}
\providecommand{\calT}{\mathcal{T}}

\section{Problem Definition}\label{sec:problem}

A new tool entering a fab passes through physical installation, supplier
qualification, and company qualification before it can process production
wafers~\citep{ieee2013semircpsp}. During a capacity ramp-up many tools go
through this sequence at the same time, and tools of the same type follow
essentially the same process network. The tools share contracted trades,
supplier engineers, company qualification engineers, and metrology equipment,
and they consume test wafers, specialty chemicals, and spare parts from a
common stock that has to be replenished.

We build the model in three steps. Section~\ref{sec:pd-mrcpsp} states the
multi-mode resource-constrained project scheduling problem (MRCPSP).
Section~\ref{sec:pd-cheng} adds the time-varying renewable capacities studied
by \citet{cheng2015splitting}. Section~\ref{sec:pd-ours} extends the result to
many projects whose non-renewable resources are replenished by orders with
lead times.

% ------------------------------------------------------------------
\subsection{The multi-mode RCPSP}\label{sec:pd-mrcpsp}

A project consists of tasks $\calJ$ with precedence arcs $E \subseteq \calJ
\times \calJ$. Task $j$ is executed in one mode $q \in \calQ_j$, which fixes
its processing time $p_{jq}$, its per-period requirement $k_{jqr}$ of each
renewable resource $r \in \calR$, and its total requirement $a_{jqm}$ of each
non-renewable resource $m \in \calN$. Renewable resource $r$ has capacity
$K_r$ in every period; non-renewable resource $m$ has a budget $N_m$ for the
whole project. With periods $\calT = \{0, \dots, H-1\}$, admissible start
periods $\calT_{jq} = \{0, \dots, H - p_{jq}\}$, and binary variables
$y_{jqt}$ equal to one if task $j$ starts in period $t$ in mode $q$, the
time-indexed formulation~\citep{pritsker1969multiproject,talbot1982resource} is
\begin{align}
    \text{(M1)} \qquad \min \quad & C_{\max} \label{eq:m1-obj} \\
    \text{s.t.} \quad
    & \sum_{q \in \calQ_j} \sum_{t \in \calT_{jq}} y_{jqt} = 1
      && \forall j \in \calJ \label{eq:m1-assign} \\
    & S_j \ge C_i
      && \forall (i,j) \in E \label{eq:m1-prec} \\
    & C_{\max} \ge C_j
      && \forall j \in \calJ \label{eq:m1-cmax} \\
    & \sum_{j \in \calJ} \sum_{q \in \calQ_j} k_{jqr}
      \sum_{\tau \in \calT_{jq} :\, \tau \le t < \tau + p_{jq}} y_{jq\tau} \le K_r
      && \forall r \in \calR,\ t \in \calT \label{eq:m1-res} \\
    & \sum_{j \in \calJ} \sum_{q \in \calQ_j} \sum_{t \in \calT_{jq}} a_{jqm}\, y_{jqt} \le N_m
      && \forall m \in \calN \label{eq:m1-nonren} \\
    & y_{jqt} \in \{0,1\}, \label{eq:m1-dom}
\end{align}
where
$S_j = \sum_{q}\sum_{t \in \calT_{jq}} t\, y_{jqt}$ and
$C_j = \sum_{q}\sum_{t \in \calT_{jq}} (t + p_{jq})\, y_{jqt}$
denote the start and completion time of task $j$. The budget
constraint~\eqref{eq:m1-nonren} is the only place non-renewable resources
appear: it limits how much is consumed in total and says nothing about when.
In the tool ramp-up setting it restricts mode choice, for example by ruling
out assigning the wafer-hungry standard qualification mode to every tool.

% ------------------------------------------------------------------
\subsection{Time-varying capacity and activity splitting}\label{sec:pd-cheng}

\citet{cheng2015splitting} generalize (M1) in two directions that matter for
Install/Qual. First, renewable capacity varies over time, and resources may be
unavailable altogether during vacations. In the time-indexed model this
replaces~\eqref{eq:m1-res} by
\begin{equation}\label{eq:m2-res}
    \sum_{j \in \calJ} \sum_{q \in \calQ_j} k_{jqr}
    \sum_{\tau \in \calT_{jq} :\, \tau \le t < \tau + p_{jq}} y_{jq\tau} \le K_{rt}
    \qquad \forall r \in \calR,\ t \in \calT.
\end{equation}
Second, they allow \emph{non-preemptive activity splitting}: a task may be
interrupted when a resource it needs becomes unavailable and resumed when the
resource returns, but it may not be interrupted at will while its resources
are available. Each activity also has a ready time and a due date, both
imposed as hard constraints. They show that the MRCPSP with non-preemptive
splitting generalizes the RCPSP with calendars, and they compare optimal
makespans with and without splitting and under full preemption using a
modified precedence-tree branch-and-bound on instances of 12 and 16
activities. Their model keeps a budget constraint of the form~\eqref{eq:m1-nonren}
for non-renewable resources (consumed in every period an activity is
processed) and minimizes the makespan of a single project.
% Checked against the paper (C&OR 53, 2015): eqs. (1)-(14) use processing
% indicators x_jt^m and mode indicators y_j^m; (7)-(8) ready times and due
% dates; (9) U_kt time-varying renewable limits; (10) non-renewable total over
% all processing periods. Some 16-activity instances were not solved in 1 h.

We adopt~\eqref{eq:m2-res} and do not allow splitting. Time-varying capacity
is what makes supplier engineers realistic: a supplier assigns field engineers
to a fab for limited periods, so a mode that needs several of them can only
run while enough are on site, and the attractive mode of a task depends on its
start time. Without splitting, a task cannot start if the capacity it needs is
missing at any point of its execution.

% ------------------------------------------------------------------
\subsection{Multi-project ramp-up with replenishable resources}\label{sec:pd-ours}

We extend the model of Section~\ref{sec:pd-cheng} in three ways.

\paragraph{(a) Many projects.}
Each tool is a project $p \in \calP$ with tasks $\calJ_p$ and arcs $E_p$;
$\calJ = \bigcup_p \calJ_p$ and $p(j)$ is the project of task $j$. Project $p$
has a release date $r_p$ (tool move-in), a due date $d_p$ (ramp target), and a
deadline $\bar d_p$, so that
$\calT_{jq} = \{ r_{p(j)}, \dots, \bar d_{p(j)} - p_{jq} \}$.
Projects are coupled only through the shared resources.

\paragraph{(b) Replenishable non-renewable resources.}
The budget $N_m$ is replaced by a stock. Material $m \in \calM$ starts with
inventory $I^0_m$; an order of $x_{mt}$ units placed in period $t$ arrives in
period $t + L_m$; each order costs $f_m$ plus $c_m$ per unit, and each unit
held at the end of a period costs $h_m$. A task consumes its materials when it
starts, and it can start only if they are on hand. With demand
\begin{equation}\label{eq:m3-demand}
    D_{mt} = \sum_{j \in \calJ} \sum_{q \in \calQ_j :\, t \in \calT_{jq}} a_{jqm}\, y_{jqt},
\end{equation}
the budget constraint~\eqref{eq:m1-nonren} becomes an inventory balance
without backlogging.

\paragraph{(c) Cost objective.}
The makespan is replaced by total cost, all in money. Project $p$ incurs a
cost of delay $w_p$ per period beyond its due date; in a fab, $w_p$ is the
value of the output the qualified tool would have produced, so bottleneck
tool types carry a larger $w_p$. Modes carry direct costs $\gamma_{jq}$.

\paragraph{Formulation.}
With order indicators $\delta_{mt}$, end-of-period inventories
$\mathrm{inv}_{mt}$, tardiness $T_p$, and
$U_m = \sum_{j} \max_{q} a_{jqm}$,
\begin{align}
    \text{(M3)} \qquad \min \quad
    & \sum_{p \in \calP} w_p T_p
      + \sum_{j \in \calJ} \sum_{q \in \calQ_j} \sum_{t \in \calT_{jq}} \gamma_{jq}\, y_{jqt}
      + \sum_{m \in \calM} \sum_{t \in \calT} \bigl( c_m x_{mt} + f_m \delta_{mt} + h_m\, \mathrm{inv}_{mt} \bigr)
      \label{eq:m3-obj} \\
    \text{s.t.} \quad
    & \text{\eqref{eq:m1-assign}},\ \text{\eqref{eq:m2-res}} \nonumber \\
    & S_j \ge C_i
      && \forall p \in \calP,\ (i,j) \in E_p \label{eq:m3-prec} \\
    & T_p \ge C_j - d_p
      && \forall p \in \calP,\ j \in \calJ_p \label{eq:m3-tard} \\
    & \mathrm{inv}_{mt} = \mathrm{inv}_{m,t-1} + x_{m,t-L_m} - D_{mt}
      && \forall m \in \calM,\ t \in \calT \label{eq:m3-inv} \\
    & x_{mt} \le U_m\, \delta_{mt}
      && \forall m \in \calM,\ t \in \calT \label{eq:m3-setup} \\
    & \mathrm{inv}_{m,-1} = I^0_m; \quad x_{mt} = 0 \ \text{ if } t < 0 \text{ or } t > H - 1 - L_m
      && \forall m \in \calM \label{eq:m3-init} \\
    & y_{jqt},\, \delta_{mt} \in \{0,1\}, \quad x_{mt} \in \mathbb{Z}_{\ge 0}, \quad
      \mathrm{inv}_{mt},\, T_p \ge 0. \label{eq:m3-dom}
\end{align}
Constraint~\eqref{eq:m2-res} now sums over the tasks of all projects and is
where the tools compete for engineers, trades, and metrology.
Constraints~\eqref{eq:m3-inv}--\eqref{eq:m3-init} describe the shared stock:
demand in periods $t < L_m$ can only be met from $I^0_m$, and later demand
from $I^0_m$ or from orders placed at least $L_m$ periods before. The
deadlines $\bar d_p$ bound each task's start window and keep the number of
variables proportional to the window rather than to $H$.

Table~\ref{tab:pd-compare} summarizes how the three models and the closest
multi-project model with material ordering~\citep{hu2025supplier} differ.

\begin{table}[t]
\centering
\caption{Model features. (M1): MRCPSP; \citet{cheng2015splitting}; \citet{hu2025supplier}; (M3): this paper.}
\label{tab:pd-compare}
\small
\begin{tabular}{@{}lcccc@{}}
\toprule
Feature & (M1) & Cheng et al. & Hu et al. & (M3) \\
\midrule
Multiple modes                          & \checkmark & \checkmark &            & \checkmark \\
Time-varying renewable capacity         &            & \checkmark &            & \checkmark \\
Non-preemptive activity splitting       &            & \checkmark &            &            \\
Activity ready times and due dates      &            & \checkmark &            &            \\
Multiple projects                       &            &            & \checkmark & \checkmark \\
Renewable resources shared across projects &         &            &            & \checkmark \\
Non-renewable: total budget             & \checkmark & \checkmark &            &            \\
Non-renewable: timed stock, orders, lead times & &            & \checkmark & \checkmark \\
Material stock shared across projects   &            &            &            & \checkmark \\
Time-varying supplier (material) capacity &          &            & \checkmark &            \\
Supplier selection                      &            &            & \checkmark &            \\
Objective                               & makespan   & makespan   & cost       & cost \\
Delay term                              &            &            & project duration & tardiness vs.\ due date \\
Solution method                         &            & B\&B       & heuristic  & B\&P \\
\bottomrule
\end{tabular}
% Hu et al. column checked against C&OR 188 (2026) 107362, eqs. (1)-(12):
% single mode (x_jt^p, fixed D_pj); renewable R_pk and inventory I_pmt are per
% project ("local resources, precluding inter-project transfer"); projects are
% coupled only through supplier capacity C_st; materials consumed every period
% of processing; objective w_p * completion + ordering + purchase + holding.
\end{table}

\paragraph{Relation between the budget and the stock.}
If no order can arrive within the horizon ($L_m \ge H$) and $h_m = 0$, every
$x_{mt}$ is zero and inventory only decreases, so~\eqref{eq:m3-inv} with
$\mathrm{inv}_{mt} \ge 0$ is equivalent to
$\sum_{j}\sum_{q}\sum_{t} a_{jqm}\, y_{jqt} \le I^0_m$, which
is~\eqref{eq:m1-nonren} with $N_m = I^0_m$. The budget of (M1) is therefore
the special case of (M3) in which replenishment is impossible. Once orders are
possible, the consumption time matters: the same total usage can be feasible
or infeasible, cheap or expensive, depending on when the tasks start.

\paragraph{Procurement for a fixed schedule.}
When $y$ is fixed, the demands $D_{mt}$ are data and the rest of (M3)
separates by material into uncapacitated single-item lot-sizing problems with
a lead time and initial stock, each solved exactly by the algorithm of
\citet{wagner1958dynamic}. The difficulty lies in the coupling: the schedule
determines when materials are needed, and the cost of meeting those needs
determines which schedule is best.

\paragraph{Assumptions.}
Durations, requirements, capacities, and lead times are deterministic. Tasks
are non-preemptive and not split. Orders are not limited by supplier capacity.

\paragraph{Size.}
(M3) has $\sum_{j}\sum_{q} |\calT_{jq}|$ binary scheduling variables and
$O(|\calM| H)$ procurement variables. Both grow with the number of tools, and
the monolithic model becomes hard to solve well before realistic ramp-up
sizes (Section~\ref{sec:experiments}). Section~\ref{sec:method} decomposes it
by project.
